\dykband{EXERCISES IN SWAFL}
{\small Open SWAFL (desktop or online) with the benchmark project. The chapter of the \href{https://github.com/halilunaziru73-creator/Coupled-Water-Fluxes-Pore-Scale-Dynamics-Precision-Irrigation-Mulching-Feedback-Loops/raw/main/Software/v1.1.0/SWAFL_Desktop/docs/manual_v1.1.0/SWAFL_Manual_v1.1.0.pdf}{SWAFL manual} that explains each task is given in brackets; record your values in a table and compare them with what is expected (\emph{Check}).}\par\vspace{3pt}
\begin{exlist}
\item {\color{khblue}[Manual, Chapter 6]} Open \emph{Sensor and trigger}. For the sand, set the trigger to 150, 300 and 600 cm and run \emph{All 20 years}. Record the irrigation and the yield. At which trigger does the irrigation stop?\par\noindent{\small\itshape Check: the sand is never irrigated above its reach of about 200 cm.}
\item {\color{khblue}[Manual, Chapter 6]} In the same layer, drag the purple trigger marker on the loam from 100 to 1000 cm and watch the next irrigation date and the depletion. Relate the interval to the readily available water.\par\noindent{\small\itshape Check: the interval grows with the trigger until the crop is stressed.}
\item {\color{khblue}[Manual, Chapter 6]} Open \emph{Trigger law} and read the depletion of the root zone at the trigger for three soils. Check that a single depletion fraction near 0.45 to 0.5 describes them.\par\noindent{\small\itshape Check: the triggers of all soils collapse onto one depletion fraction.}
\item {\color{khblue}[Manual, Chapter 18]} Use the search for the most efficient inputs on the drip layer of the loam with the trigger free. Which trigger does the search choose, and how does it compare with the reach?\par\noindent{\small\itshape Check: the chosen trigger lies well inside the reach.}
\end{exlist}
